\section{Audit of the Current Proposal}
\label{sec:audit}

Two versions of the slide deck were supplied: Deck~A (DenseNet121 vs.\ a ``Sequential'' CNN, image classification) and Deck~B (YOLOv8 vs.\ Faster R-CNN with Grad-CAM, object detection on leaf, fruit and pest images). Both are \emph{technically executable}; the problems are in experimental protocol, internal consistency, and novelty. Table~\ref{tab:audit} lists what must be fixed before any of it can be written up.

\begin{longtable}{@{}L{0.5cm}L{5.2cm}L{5.4cm}L{5.0cm}@{}}
\caption{Issues found in Deck~A and Deck~B, with the fix adopted in this plan.}\label{tab:audit}\\
\toprule
\# & Issue & Why it matters & Fix in this plan \\
\midrule
\endfirsthead
\toprule
\# & Issue & Why it matters & Fix in this plan \\
\midrule
\endhead
\bottomrule
\endfoot
\multicolumn{4}{@{}l}{\textbf{Deck A -- DenseNet121 classification}}\\
\midrule
A1 & Three different accuracies for the same DenseNet121 run: 99.6\% (conclusion), 99.6/99.2\% train/val (results), 96.2/94\% (table). & Reviewers reject papers whose numbers disagree with themselves. & One results table, generated by a script from saved predictions; every number traceable to a run ID and seed. \\
A2 & Dataset described as ``high-resolution RGB photos captured in natural environments''. The 10\,000-image, 10-class set is a Kaggle re-split of PlantVillage: 256$\times$256 lab images on uniform backgrounds. & Mis-describing the data hides the biggest weakness (lab-only images do not transfer to the field; see \S\ref{sec:sota}). & Describe each dataset by provenance, acquisition conditions and known biases (\S\ref{sec:datasets}); add field datasets. \\
A3 & Learning rates written as ``1e-4 (0.00001)'' and ``1e-3 (0.0001)''. & Arithmetic errors in hyper-parameter tables undermine trust. & Hyper-parameters stored in a versioned YAML; table generated from it. \\
A4 & The ``Sequential'' baseline (85\%) is a straw man; the claim ``outperformed state-of-the-art'' is made without comparing to any published number. & PlantVillage-tomato classification is saturated at 99--99.9\% in the literature; 96--99\% is not a contribution. & Compare against the full baseline set in \S\ref{sec:baselines}, on identical splits, with mean$\pm$std. \\
A5 & ``When the gradient vanishing issue is fixed'' -- ReLU inside DenseNet is the stock architecture. & Presented as a contribution; it is not. & Remove. Novelty is defined in \S\ref{sec:novelty}. \\
A6 & 80/20 train/test with ``validation accuracy'' reported; no separate validation split; no repeated runs. & Model selection on the test set inflates results; a single run has no error bar. & 70/15/15 stratified split with deduplication; 5 seeds (or 5-fold CV); McNemar / paired test between models. \\
A7 & Precision 0.98, recall 0.99, F1 0.99 with accuracy 99.6\% -- inconsistent (macro-F1 cannot exceed both P and R). & Metric definitions unclear (micro vs macro). & Report macro-P/R/F1, kappa, MCC, per-class confusion matrix (\S\ref{sec:metrics}). \\
\midrule
\multicolumn{4}{@{}l}{\textbf{Deck B -- YOLOv8 / Faster R-CNN + Grad-CAM}}\\
\midrule
B1 & ``Accuracy 94\%'' reported for a detector; fruit accuracy is 87\% in the table and 89\% in the text; Grad-CAM IoU 0.40 in the table, 0.39 in the text; ``NSR'' never defined. & Accuracy is not a detection metric unless defined (image-level? box-level at which IoU?); inconsistent values. & Report COCO-style mAP@0.5, mAP@0.5:0.95, AP$_s$/AP$_m$/AP$_l$, P/R/F1 at a fixed confidence; define every XAI metric with a formula (\S\ref{sec:metrics}). \\
B2 & Grad-CAM applied to YOLOv8 without stating the target layer, the target score (class logit? objectness? which box?) or how gradients pass NMS. & Grad-CAM needs a scalar target; for one-stage detectors the standard practice is EigenCAM / Grad-CAM on the neck output for a chosen box. Unspecified = irreproducible. & Specify: EigenCAM and Grad-CAM++ on the P3/P4 neck features, target = class score of the selected box before NMS; report pointing game and energy-based pointing game against lesion masks. \\
B3 & ``Avg.\ Grad-CAM IoU 0.51'' requires ground-truth lesion masks; only boxes exist. & If IoU was computed against boxes, it is a box-IoU of a thresholded heatmap, not a lesion-localisation score. & Annotate a small lesion-mask subset (200--300 images) for XAI evaluation, or use the pixel-mask datasets in \S\ref{sec:datasets}. \\
B4 & Pipeline text describes a ``classification network with self-attention, multi-head attention and transfer learning'' after YOLOv8, but no such network appears in results. & Claimed component not evaluated. & Replace with the actual custom model (\S\ref{sec:model}); every named component gets an ablation row. \\
B5 & Faster R-CNN (2015) is the only comparator. & YOLOv8 $>$ Faster R-CNN in speed and usually mAP is a known result; not a finding. & Baselines: YOLOv5/8/10/11/12, RT-DETR, plus lightweight classifiers (\S\ref{sec:baselines}). \\
B6 & Fruit dataset is ripe/unripe (ripeness), yet results are described as ``fruit disease detection''. & Task mismatch. & Either annotate fruit disorders (blossom-end rot, cracking, sunscald) or present the fruit branch as ripeness/maturity, not disease. \\
B7 & Custom North-Gujarat set (3\,153 ``mixture'' images) has no described class list, annotation protocol, camera, season, or inter-annotator agreement. & This set is the strongest novelty asset in the whole project; undocumented it cannot be published or reused. & Dataset card + annotation protocol (\S\ref{sec:custom_dataset}); release on Zenodo/Mendeley with DOI. \\
B8 & Pre-processing says ``one-hot encoded vectors for multi-class classification'' for a detector; 640$\times$640 for classification too. & Copy--paste from Deck~A; detectors use box labels; classifiers normally use 224--384 px. & Task-specific pre-processing (\S\ref{sec:pipeline}). \\
B9 & Research-gap slide lists severity, multi-disease-per-leaf, field robustness, temporal tracking, explainability -- but the proposed method only addresses explainability (qualitatively). & Gap and method are disconnected. & Each claimed gap maps to a measurable experiment (Table~\ref{tab:gap_map}). \\
B10 & Bibliography: duplicate [21][21]; numbering differs between decks; several ``research observation'' cells contradict the paper described (e.g.\ Reis~2024 listed as KNN/SVM but described as lightweight CNN + meta-learning; Nawaz~2023 ``highly effective'' in Deck~A, ``has not proven effective'' in Deck~B). & Credibility. & BibTeX with DOIs; observation column rewritten from the actual abstracts. \\
\end{longtable}

\paragraph{Verdict on feasibility.}
\begin{itemize}
  \item \textbf{Doable now:} YOLOv8/Faster R-CNN detection with Grad-CAM (already run), DenseNet121 classification (already run).
  \item \textbf{Doable with moderate effort (2--4 months on one GPU):} the custom dual-task model of \S\ref{sec:model}, the cross-domain protocol, SSL pre-training, quantitative XAI.
  \item \textbf{Not doable as stated:} ``fruit disease'' without disease-annotated fruit; ``severity'' without lesion masks or expert severity scores; ``temporal progression'' without a time-series dataset. These are dropped or scoped down explicitly in \S\ref{sec:experiments}.
  \item \textbf{Where the novelty actually is:} (i) a single block family that instantiates both a classifier and a detector, trained under one recipe and evaluated under one protocol; (ii) a lab$\rightarrow$field generalisation study across five public datasets plus the North-Gujarat field set; (iii) self-supervised pre-training on pooled unlabelled tomato images with label-efficiency curves; (iv) healthy-only anomaly detection for \emph{unknown} disease flagging; (v) quantitative, not qualitative, explainability. No single one of these is new; the combination on tomato, with a released field dataset, is publishable.
\end{itemize}
